\centerline{\bf \Huge Правильный перебор}

\problem{\textbf{1}} Начальник обратил внимание, что дата проведения десятого дня IV Летних математических сборов, записанная восемью цифрами (30.07.2026) обладает интересной особенностью: сумма первых четырёх цифр равна сумме цифр года. А какие ещё даты в 2026 году имеют такое же свойство? Перечислите ВСЕ варианты!\\

\problem{\textbf{2}} Кузнечик прыгает вдоль прямой вперёд на 80 см или назад на 50 см. Может ли он менее чем за 7 прыжков удалиться от начальной точки ровно на 1 м 70 см?\\

\problem{\textbf{3}} Перечислите все четвёрки натуральных чисел, дающих в сумме
15. 

Например: 1+1+2+11=15.\\


\problem{\textbf{4}} У 2009 года есть такое свойство: меняя местами цифры числа 2009, нельзя получить меньшее четырехзначное число (с нуля числа не начинаются). В каком году это свойство впервые повторится снова?\\

\problem{\textbf{5}}  Напишите в строчку первые 10 простых чисел. Как вычеркнуть 6 цифр, чтобы получилось наибольшее возможное число?\\

\problem{\textbf{6}} В январе некоторого года было 4 понедельника и 4 пятницы.
Каким днем недели могло быть 20-е число этого месяца?\\


\problem{\textbf{7}}  На конкурсе «А ну-ка, чудища!» стоят в ряд 15 драконов. У соседей число голов отличается на 1. Если у дракона больше голов, чем у обоих его
соседей, его считают хитрым, если меньше, чем у обоих соседей, — сильным, остальных (в том
числе стоящих с краю) считают обычными. В ряду есть ровно четыре хитрых дракона — с 4,
6, 7 и 7 головами и ровно три сильных — с 3, 3 и 6 головами. У первого и последнего драконов
голов поровну. Приведите хотя бы два примера того, как такое могло быть.

